\documentclass[12pt,letterpaper]{article}
\usepackage[utf8]{inputenc}
\usepackage[T1]{fontenc} 
\usepackage[margin=2.5cm]{geometry} 
\usepackage{graphicx}
\usepackage{float}
\usepackage{tabularx}
\usepackage{makecell}

\title{Referencia Técnica}
\author{Fabian Ochoa}
\date{May 2026}

\begin{document}

\maketitle

\section{Rendimiento de Computadores}

\renewcommand{\arraystretch}{1.5}

\subsection{Prestaciones y Tiempo}
\textbf{¿Para qué sirve?:} Define la relación fundamental entre el rendimiento absoluto de un computador y el tiempo que tarda en ejecutar una tarea.

\begin{table}[H]
    \centering
    \begin{tabularx}{\textwidth}{|l|X|}\hline
         \textbf{FÓRMULA} & \textbf{EXPLICACIÓN} \\\hline
         \texttt{Prestaciones(X) = 1 / Tiempo(X)} & Las prestaciones son inversamente proporcionales al tiempo; para maximizar el rendimiento se debe minimizar el tiempo. \\ \hline
    \end{tabularx}    
\end{table}

\subsection{Regla de Comparación entre Máquinas}
\textbf{¿Para qué sirve?:} Establece la regla lógica de comparación entre dos máquinas para determinar cuál es más rápida.

\begin{table}[H]
    \centering
    \begin{tabularx}{\textwidth}{|l|X|}\hline
         \textbf{FÓRMULA} & \textbf{EXPLICACIÓN} \\\hline
         \makecell[l]{\texttt{Prestaciones(X) > Prestaciones(Y)} \\ \texttt{implica} \\ \texttt{Tiempo(Y) > Tiempo(X)}} & Si una máquina X es más rápida que Y, significa que la máquina Y requiere más tiempo para completar la misma tarea. \\ \hline
    \end{tabularx}    
\end{table}

\subsection{Cálculo de Velocidad Relativa usando Prestaciones}
\textbf{¿Para qué sirve?:} Calcula cuántas veces (n) es más rápida la máquina X que la Y utilizando directamente la relación entre sus prestaciones.

\begin{table}[H]
    \centering
    \begin{tabularx}{\textwidth}{|l|X|}\hline
         \textbf{FÓRMULA} & \textbf{EXPLICACIÓN} \\\hline
         \texttt{n = Prestaciones(X) / Prestaciones(Y)} & Se utiliza una proporción directa entre sus capacidades de rendimiento. \\ \hline
    \end{tabularx}    
\end{table}

\subsection{Cálculo de Velocidad Relativa usando Tiempos}
\textbf{¿Para qué sirve?:} Calcula cuántas veces (n) es más rápida la máquina X que la Y utilizando la inversa de los tiempos de ejecución.

\begin{table}[H]
    \centering
    \begin{tabularx}{\textwidth}{|l|X|}\hline
         \textbf{FÓRMULA} & \textbf{EXPLICACIÓN} \\\hline
         \texttt{n = Tiempo(Y) / Tiempo(X)} & Dado que medir prestaciones es complejo, se usa la inversa del tiempo de la máquina más lenta respecto a la más rápida. \\ \hline
    \end{tabularx}    
\end{table}

\end{document}